\documentclass[conference]{IEEEtran}
\IEEEoverridecommandlockouts

\usepackage{cite}
\usepackage{amsmath,amssymb,amsfonts}
\usepackage{algorithmic}
\usepackage{algorithm}
\usepackage{graphicx}
\usepackage{textcomp}
\usepackage{xcolor}
\usepackage{booktabs}

\def\BibTeX{{\rm B\kern-.05em{\sc i\kern-.025em b}\kern-.08em
    T\kern-.1667em\lower.7ex\hbox{E}\kern-.125emX}}

\begin{document}

\title{Automated Schedule Generation for Lithography Machine
Upgrades: A Comparative Study of Priority-Rule Greedy, CP-SAT, and a
Domain-Enhanced Greedy Heuristic%
\thanks{This thesis was prepared in partial fulfilment of the requirements
for the Degree of Bachelor of Computer Science, Maastricht University.
Supervisor(s): Guangzhi Tang, Chang Sun.}
}

\author{\IEEEauthorblockN{Xuhan Zhuang}
\IEEEauthorblockA{\textit{Department of Advanced Computing Sciences} \\
\textit{Faculty of Science and Engineering}\\
\textit{Maastricht University}\\
Maastricht, The Netherlands}
}

\maketitle

% ============================================================
\begin{abstract}
Planning a downtime upgrade for a complex lithography machine
requires scheduling tens to hundreds of interdependent tasks under
tight resource and physical constraints. We formalise one such
problem, arising at ASML, as a Resource-Constrained Project
Scheduling Problem (RCPSP) with five domain-specific extensions:
single-day non-preemption, time-varying staffing capacity, spatial
exclusivity, tool exclusivity, and machine-state compatibility. Three
candidates are evaluated on four real industrial instances: a
priority-rule greedy constructor with Most Total Successors (MTS) and
Shortest Processing Time (SPT) variants, an exact CP-SAT solver, and
a domain-enhanced greedy heuristic. The recommended heuristic
configuration attains a 6.1\% mean optimality gap to the CP-SAT
incumbent across the four instances, well below the pre-registered
15\% threshold and a 37\% relative reduction over the best classical
priority rule, all at sub-second cost. CP-SAT proves optimality in
seconds on three of the four instances but exhausts a 600-second
budget without closing the gap on the fourth. A budget sweep reveals
complementary regimes: in the sub-second window relevant for
interactive scenario iteration, CP-SAT returns no feasible solution
on three instances and the heuristic is the only viable choice; from
one second of compute upward, CP-SAT is never worse than the
heuristic on any instance.
\end{abstract}

\begin{IEEEkeywords}
resource-constrained project scheduling, constraint programming,
CP-SAT, domain-specific heuristics, makespan minimisation,
lithography manufacturing
\end{IEEEkeywords}

% ============================================================
\section{Introduction}

ASML is a manufacturer of semiconductor lithography machines.
Industrial engineers there currently schedule the machine upgrade
process manually. The resulting schedule is referred to internally as
a high-level sequence. The process is slow, error-prone, and does not
scale as machines grow more complex. Schedule quality determines how
long the machine is out of service, which in semiconductor production
is the dominant economic cost. This paper investigates algorithms to
replace this manual process.

The core difficulty is that, beyond classical precedence, the
scheduler must simultaneously satisfy single-day non-preemption,
staffing capacity across working hours, exclusive access to physical
workspace zones, exclusive access to specialised tools, and
machine-state compatibility (each task declares a value on a small set
of states, and concurrent tasks must agree on every state they both
declare). This combination positions the problem as a
Resource-Constrained Project Scheduling Problem
(RCPSP)~\cite{brucker1999rcpsp} with five domain-specific extensions.

RCPSP is NP-hard, and decades of research have produced two broad
solution families: exact solvers that prove optimality but whose cost
grows sharply with instance size and constraint interaction, and
faster heuristic constructions that trade optimality guarantees for
speed. Section~\ref{sec:related} reviews both families and positions
the present formulation against them.

Two industrial considerations motivate studying heuristic alternatives
alongside the exact solver. First, scalability is uncertain: as ASML
machines grow more complex, it is not known whether an exact solver
scales well across instances of increasing size under all five
extension constraints. Second, speed matters in production: an engineer
often sweeps many planning variants within a single session, so a fast
heuristic with a small loss in solution quality can be preferable to a
slow exact solver even when the latter returns the optimum.

\subsection{Research question and hypotheses}
\label{sec:rqs}

This thesis addresses the following research question through a
comparative algorithmic study.

\smallskip
\noindent\textbf{RQ.} \emph{To what extent can a domain-enhanced
greedy heuristic close the optimality gap between classical
priority-rule constructive heuristics and the CP-SAT
constraint-programming solver in Google OR-Tools across real
industrial instances, at sub-second computational cost?}
\smallskip

Three candidates are evaluated: (A) a baseline greedy constructor
using standard RCPSP priority rules, (B) Google OR-Tools CP-SAT as an
exact benchmark, and (C) a domain-enhanced greedy heuristic (DEH) whose
composite scoring function encodes problem-specific constraints
directly into task selection. Two falsifiable hypotheses about the
heuristic are pre-registered, then evaluated on four real industrial
instances drawn from distinct machine modules.

\begin{itemize}
  \item \textbf{H1.} Over the four evaluated instances, the mean
  optimality gap of the best-configured domain-enhanced greedy
  heuristic against the CP-SAT solver is strictly smaller than
  that of both classical priority rules.
  \item \textbf{H2.} Over the four evaluated instances, the mean
  optimality gap of the best-configured domain-enhanced greedy
  heuristic against the CP-SAT solver does not exceed
  15\%.\footnote{Kolisch and Hartmann~\cite{kolisch2006experimental}
  report 10--14\% mean gaps for the best classical priority rules on
  PSPLIB; 15\% sits just above that band.}
\end{itemize}

This study makes four contributions. First, it formalises the ASML
high-level-sequence problem as an RCPSP with five domain-specific
extensions, including a machine-state compatibility constraint not
present in standard formulations. Second, it evaluates three
candidates (two classical priority rules, an exact CP-SAT model, and a
domain-enhanced greedy heuristic) under identical inputs on four real
industrial instances. Third, it characterises the speed--quality
trade-off through a CP-SAT time-budget sweep, identifying the
sub-second regime in which the heuristic is the only candidate that
returns a feasible schedule. Fourth, it reports transparently which
signals drive the result: a critical-path ordering closes most of the
gap, while the additional domain-specific terms prove at best
conditional on the current instances.

\subsection{Paper structure}

Section~\ref{sec:related} reviews background and related work.
Section~\ref{sec:methods} formalises the problem and describes the
three algorithmic approaches. Section~\ref{sec:experiments} describes
the experimental design and metrics. Section~\ref{sec:results}
presents results across the four available instances.
Section~\ref{sec:discussion} interprets the findings and discusses
limitations and future work. Section~\ref{sec:conclusions} concludes.

% ============================================================
\section{Background and Related Work}
\label{sec:related}

\subsection{Resource-Constrained Project Scheduling}

In the classical RCPSP a set of tasks with fixed durations must be
scheduled subject to precedence constraints and a limited supply of
renewable resources, so as to minimise the makespan. The problem
generalises job-shop scheduling and is NP-hard in the strong
sense~\cite{blazewicz1983scheduling}. Standard notation, a
classification scheme, and a catalogue of models and methods are given
by Brucker et al.~\cite{brucker1999rcpsp}, with a book-length treatment
in the research handbook of Demeulemeester and
Herroelen~\cite{demeulemeester2002project}. Most empirical work in the
field is benchmarked on PSPLIB~\cite{kolisch1997psplib}, a library of
generated instances with known or best-known makespans, which has
shaped how heuristics are compared. The instances studied here are real
industrial machine modules at ASML rather than instances drawn from
PSPLIB, so these library baselines do not apply here.

\subsection{Exact Methods}

Exact approaches solve the problem to proven optimality. Early work used
dedicated branch-and-bound, such as the procedure of Stinson et
al.~\cite{stinson1978multiple} for the multiple-resource case.
Mixed-integer programming and constraint programming are now the
standard exact tools. Constraint-based scheduling, surveyed by
Baptiste et al.~\cite{Baptiste2001}, models each task as an interval
variable and reasons over global constraints such as \emph{cumulative}
for renewable resources and \emph{no-overlap} for disjunctive ones.
Modern CP solvers such as IBM ILOG CP Optimizer~\cite{laborie2018ibm}
and Google OR-Tools CP-SAT~\cite{ortools} pair these constraints with
strong filtering and systematic search. The cumulative constraint in
particular depends on specialised filtering, of which the edge-finding
rule of Vil\'im~\cite{vilim2015fastcumulative} and the conditional
time-interval reasoning of Laborie and
Rogerie~\cite{laborie2008complete} are representative. Such solvers
prove optimality on small and moderate instances, but their cost grows
quickly with size and constraint interaction, which is the behaviour we
probe directly in the budget sweep of Section~\ref{sec:budget-sweep}.

\subsection{Constructive and Metaheuristic Methods}

For larger or time-critical instances, heuristics remain the practical
standard. The simplest are priority-rule methods built on a
schedule-generation scheme: tasks are placed one at a time in an order
set by a priority rule, with the serial and parallel schemes differing
in how the next task is chosen. Kolisch~\cite{kolisch1996efficient}
studies the efficiency of common priority rules, among them the
most-total-successors and shortest-processing-time rules that we adopt
as the two baselines of Candidate A. Metaheuristics such as genetic
algorithms, tabu search, and simulated annealing then search the space
of activity orderings; the experimental survey of Kolisch and
Hartmann~\cite{kolisch2006experimental} compares many of these on
PSPLIB and reports mean optimality gaps of roughly 10 to 14\% for the
best priority rules. Hartmann and Briskorn~\cite{hartmann2010survey}
survey the broader space of variants and extensions. Our heuristic
stays in this cheap constructive class rather than the metaheuristic
one, since the target is fast, lightweight scheduling rather than the
longer search a metaheuristic performs.

\subsection{Industrial Project Scheduling Applications}

Beyond benchmark studies, RCPSP models and their extensions have been
applied to industrial planning problems that share structural features
with the present one. Maintenance turnaround and shutdown projects, in
which a facility is taken offline and a large set of interdependent jobs
must be completed under crew and equipment limits before it returns to
service, have been modelled as resource-constrained project
scheduling~\cite{megow2011turnaround}. Scheduling in
semiconductor manufacturing has likewise drawn sustained attention, for
both production and equipment operations, including problems with
auxiliary resources and cluster tools~\cite{monch2011semiconductor}.
In these settings the planning objective is typically to minimise the
downtime or completion window, the direct analogue of the makespan
minimised here, which is what makes RCPSP the natural modelling frame
for the ASML upgrade scheduling problem.

\subsection{Extensions and Positioning}

Real problems rarely match the textbook RCPSP, and a large literature
studies extensions: time-varying resource availability, multiple
execution modes, generalised precedence, and non-renewable or
partially renewable resources, all catalogued
in~\cite{hartmann2010survey}. Each constraint we impose has appeared
individually in prior work. Spatial and tool exclusivity are
disjunctive (no-overlap) resources, time-varying staffing is a calendar
on a renewable resource, and machine-state compatibility resembles a
sequence-dependent setup. What we have not found in the published
literature is this particular combination treated as a single
formulation, in particular machine-state compatibility together with
time-varying staffing capacity. Because instance structure drives both
solver difficulty and heuristic quality, and because generated
benchmarks need not reflect real structure~\cite{vanhoucke2008measurement},
we evaluate on real modules chosen to span complementary difficulty
profiles rather than on a single synthetic family.

% ============================================================
\section{Methods}
\label{sec:methods}

\subsection{Problem Formalisation}

The ASML upgrade scheduling problem is a Resource-Constrained Project
Scheduling Problem (RCPSP) with five domain-specific extensions to
the classical formulation. In classical RCPSP, a set of tasks must
be scheduled subject to precedence and renewable resource
constraints, minimising total project duration (makespan). The
formulation here adds single-day non-preemption, time-varying
staffing capacity, spatial exclusivity, tool exclusivity, and
machine-state compatibility.

\subsubsection{Inputs}

The algorithm receives the available working time as an ordered list
of \emph{working days}, together with a \emph{task set} $T$ defining
the work to be scheduled. Each working day offers a fixed block of
working hours (the same on every day), staffed by an integer number of
engineers available throughout; hours outside the block have no staff
and admit no work.

Each task $i \in T$ declares an integer duration $d_i$ in hours; a
staff requirement $f_i$ (the number of engineers required throughout);
a value $p_i^{(a)}$ for each independent machine-state
\emph{variable} $a$, with $p_i^{(a)} = \emptyset$ if the task leaves
variable $a$ unset; a set of workspace zones; and at most one required
tool, either \emph{exclusive} (only one exists, so tasks requiring it
cannot overlap) or \emph{shared} (effectively unlimited, imposing no
constraint). Precedence over $T$ forms a directed acyclic
graph $(T, E)$, in which an edge $(i, j) \in E$ means task $i$ must
finish before task $j$ may start.

\subsubsection{Constraints}

Every algorithm enforces six hard constraints. A schedule is feasible
only if all six hold for all placed tasks.

\begin{enumerate}
  \item \textbf{Single-day non-preemption.} Each task executes as a
  single uninterrupted interval inside one day, entirely within the
  working hours of that day.
  \item \textbf{Time-varying staffing capacity.} At every working
  hour, the sum of staff requirements across concurrent tasks
  must not exceed that day's staff capacity (one staff unit is one
  engineer working that day).
  \item \textbf{Spatial exclusivity.} Each task occupies a (possibly
  empty) set of workspace zones. Two tasks whose zone sets intersect
  may not execute concurrently.
  \item \textbf{Tool exclusivity.} No two tasks requiring the same
  exclusive tool may overlap. Tasks requiring a
  shared tool or no tool place no such constraint.
  \item \textbf{Precedence.} If task $i$ precedes task $j$, then $j$
  may not start until $i$ completes.
  \item \textbf{Machine-state compatibility.} Each task fixes some
  state variables to required values and leaves the rest unset (an
  unset variable imposes no constraint). Two concurrent tasks must
  have matching values on every state variable that both fix; if
  either leaves a variable unset, that variable raises no conflict.
  Any mismatch forbids overlap.
\end{enumerate}

These constraints model reconfiguration between sequential tasks as
instantaneous; switching costs across non-overlapping states
are out of scope for the present formulation.

\subsubsection{Objective}

Minimise the makespan in calendar days, defined as the day index of
the latest-completing task (the number of calendar days the machine
is out of service). This is standard makespan minimisation for
RCPSP, with the working-day discretisation matching how downtime is
reported in the field.

\subsection{Candidate Algorithms}

All three algorithms solve the same scheduling task, with identical
inputs and identical output structures. The two greedy schedulers
(Candidates A and C) break ties by ascending task identifier, giving
deterministic output.

\subsubsection{Candidate A: Baseline Priority-Rule Greedy}

A priority-rule greedy scheduler maintains the set of schedulable
tasks (those whose predecessors have all been placed), selects one
by a priority rule $\pi$, assigns its earliest feasible start time,
and advances. Two priority rules are evaluated: \textbf{Most Total Successors
(MTS)}, which prioritises tasks with the largest count of direct
and indirect successors, and \textbf{Shortest Processing Time
(SPT)}, which prioritises tasks with the shortest duration. With
$\pi(i)$ as a generic priority, MTS sets $\pi(i)$ to the size of
the successor closure of $i$ and SPT sets $\pi(i) = -d_i$, so that
$\arg\max$ in Algorithm~\ref{alg:greedy} selects the MTS-largest or
SPT-shortest task as required. The
horizon is assumed long enough that every task admits a feasible
placement; the safety-net horizon used throughout is described in
Section~\ref{sec:shiftconfig}. The full procedure is given in
Algorithm~\ref{alg:greedy}.

\begin{algorithm}[t]
\caption{Priority-rule greedy scheduler (Candidate A).}
\label{alg:greedy}
\begin{algorithmic}[1]
\REQUIRE Task set $T$, precedence edges $E$, priority rule $\pi$
\STATE $P \gets \emptyset$ \COMMENT{placed tasks}
\STATE $S \gets \{i \in T : i \text{ has no predecessors}\}$
\WHILE{$S \neq \emptyset$}
  \STATE $i^{*} \gets \arg\max_{i \in S}\, \pi(i)$
  \STATE $t_{i^{*}} \gets$ earliest hourly start, in input day order, where $i^{*}$ is feasible under all six constraints
  \STATE $P \gets P \cup \{i^{*}\}$;\quad $S \gets S \setminus \{i^{*}\}$
  \STATE $S \gets S \cup \{j \notin P \cup S : \text{all predecessors of } j \text{ are in } P\}$
\ENDWHILE
\RETURN start times $\{t_i\}_{i \in P}$
\end{algorithmic}
\end{algorithm}

\subsubsection{Candidate B: Constraint Programming with CP-SAT}

The exact baseline uses Google OR-Tools CP-SAT~\cite{ortools}. Each
task is modelled as an interval variable (a CP-SAT primitive with
start, end, and duration fields) with fixed duration; its start time
may only fall where the task fits inside the working hours of a single
day. Let $C(t)$ be the available staff at time $t$, equal to the day's
staff count during working hours and zero outside them. The staffing
limit is a cumulative constraint~\cite{vilim2015fastcumulative}
enforcing $\sum_i f_i \cdot \mathbb{1}[t \in [t_i, t_i + d_i)] \leq
C(t)$ across the scheduling horizon.

Spatial and tool exclusivity become no-overlap
constraints~\cite{laborie2008complete} (which forbid any two
intervals in a group from running concurrently) over zone-grouped
and tool-grouped intervals. Precedence is a linear inequality $t_j
\geq t_i + d_i$ per edge. Machine-state constraints add a
no-overlap between any two tasks with different declared values on a
shared state variable; on each variable, tasks declaring $\emptyset$
are unconstrained on that variable only.

The objective is the makespan in calendar days. The solver runs
synchronously with a user-supplied wall-clock time limit (600 seconds
in our experiments); if it terminates inside the limit the solution
is provably optimal, otherwise the best feasible solution is returned
along with a lower bound.

\subsubsection{Candidate C: Domain-Enhanced Greedy Heuristic (DEH)}

Candidate C runs the same construction as Algorithm~\ref{alg:greedy},
replacing the priority rule $\pi$ with a composite score. Classical
rules order tasks only by graph and duration structure, so they ignore
the machine-state constraint until the feasibility check; this can
place a state-conflicting task where a compatible one would have
avoided a day-break. The score blends a critical-path remaining term
(CPR) with a machine-state affinity term (MSA) under a tunable weight
$w$, plus an optional constraint-tightness term (CST). The three terms
are defined next, and their weights are characterised empirically in
Section~\ref{sec:results}.

\paragraph{Critical Path Remaining}
Before scheduling, the longest-path distance from every task to a
terminal node of the precedence graph is computed, measured in
cumulative processing duration:
\begin{equation}
\phi_{\text{cpr}}(i) = d_i + \max_{(i,j) \in E} \phi_{\text{cpr}}(j),
\end{equation}
with $\phi_{\text{cpr}}(i) = d_i$ for terminal tasks. We use the
normalised form $\hat{\phi}_{\text{cpr}}(i) = \phi_{\text{cpr}}(i) /
\max_j \phi_{\text{cpr}}(j) \in [0,1]$ in the composite score below.
CPR is a duration-weighted analogue of MTS: it weights tasks by the
total duration of the longest path to a terminal rather than by
successor count, so long downstream paths dominate the priority
rather than fan-out alone.

\paragraph{Machine-State Affinity}
Although constraint 6 only forbids concurrent conflicts, the machine
retains a state between tasks. MSA rewards a candidate task
for declaring state values that agree with this current state.
The algorithm maintains, for each state variable $a$ independently, the value
$p_{\text{cur}}^{(a)}$ declared by the latest-starting placed task
that declared a value on $a$. If no placed task has yet declared a value on $a$,
$p_{\text{cur}}^{(a)} = \emptyset$ (``unset''). For a candidate task $i$ that sets a non-empty set of
state variables $A_i$, the MSA score is the fraction of those variables whose
declared value $p_i^{(a)}$ either matches $p_{\text{cur}}^{(a)}$ or
for which the variable is still unset:
\begin{equation}
\phi_{\text{msa}}(i) =
\frac{1}{|A_i|} \sum_{a \in A_i}
\mathbb{1}\!\left[\, p_i^{(a)} = p_{\text{cur}}^{(a)} \,\lor\,
                       p_{\text{cur}}^{(a)} = \emptyset \,\right].
\end{equation}
Tasks that set no state variables are state-neutral and score
$\phi_{\text{msa}}(i) = 1$.

\paragraph{Constraint Tightness}
CST is a static, per-task measure of placement difficulty:
\begin{equation}
\phi_{\text{cst}}(i) = \frac{1}{3} \left(
  \frac{|A_i|}{\max_j |A_j|} +
  \frac{f_i}{\max_j f_j} +
  \frac{d_i}{\max_j d_j}
\right),
\end{equation}
where each feature is normalised by the task set's own maximum (a term with a zero maximum contributes 0). The
three features are the per-task quantities that most directly drive
per-step infeasibility risk: variable-rich tasks restrict
state-compatible partners, staff-heavy tasks consume staffing
capacity, and long-duration tasks foreclose late-day slots. Tasks
heavy on all three score near 1. Its effect on ordering is not fixed a priori, so CST is an optional
tunable boost, characterised empirically (Section~\ref{sec:results}).

\paragraph{Composite Score}
The composite score is a convex combination of CPR and MSA,
optionally augmented by the CST term:
\begin{equation}
\text{score}(i) = (1 - w) \cdot \hat{\phi}_{\text{cpr}}(i)
                + w \cdot \phi_{\text{msa}}(i)
                + \gamma \cdot \phi_{\text{cst}}(i),
\label{eq:composite}
\end{equation}
where $w \in [0,1]$ controls the trade-off between critical-path
priority and state affinity, and $\gamma \geq 0$ controls the
CST boost. The convex form keeps CPR and MSA on the same $[0,1]$
scale so that $w$ has a direct interpretation as the fractional
emphasis placed on state affinity. CST is kept additive so that
$\gamma$ varies independently of $w$ (the CST ablation is in
Section~\ref{sec:cross-instance}). Scores order tasks; the
highest-scoring task is selected at each step and assigned its
earliest feasible start time.
When $w = 0$ and $\gamma = 0$, the algorithm is pure CPR; when
$w = 1$ and $\gamma = 0$, it is pure state affinity, a signal
classical priority rules cannot express.

% ============================================================
\section{Experimental Setup}
\label{sec:experiments}

\subsection{Instances}
\label{sec:instances}

The study uses real industrial instances at ASML drawn from distinct
machine modules. Each instance is supplied as a three-table dataset
(task properties, a precedence matrix, and a spatial-zone matrix)
used internally by ASML.

Four instances are evaluated, MOD\_A through MOD\_D, chosen to span
complementary regions of the design space rather than repeated draws
from a single profile. MOD\_A is the largest and most constraint-rich,
exercising every constraint type the formulation defines; MOD\_B has a
shallow DAG and no state activity, so its difficulty lies in spatial
and tool exclusivity; MOD\_C is small but precedence-dense for its
size; and MOD\_D is the most precedence-heavy with negligible state
activity, isolating dense precedence from machine-state interaction.
Table~\ref{tab:instances} summarises their structural features.

\begin{table}[ht]
\caption{Structural features of evaluated instances.
$|T|$: task count. $\sum d_i$: total processing work in hours.
$d_{\max}$, $f_{\max}$: largest task by duration and staff requirement.
$\overline{|A_i|}$: mean state variables per task. Depth: longest precedence
chain. Density: $|E|/\binom{|T|}{2}$, where $|E|$ is the precedence
edge count.
$Z$: zone count. $T_{\text{ex}}$: exclusive tool count.}
\label{tab:instances}
\centering
\renewcommand{\arraystretch}{1.15}
\setlength{\tabcolsep}{4pt}
\begin{tabular}{@{}lrrrrrrrrr@{}}
\toprule
Inst.\ & $|T|$ & $\sum d_i$ & $d_{\max}$ & $f_{\max}$ & $\overline{|A_i|}$ & Depth & Density & $Z$ & $T_{\text{ex}}$ \\
       &       & (h)        & (h)        & (staff)    &                    &       &         &     &                 \\
\midrule
MOD\_A & 100 & 458 & 10 & 5 & 0.91 & 29 & 0.043 & 27 & 6 \\
MOD\_B &  90 & 300 & 10 & 5 & 0.00 & 11 & 0.035 & 31 & 1 \\
MOD\_C &  22 &  86 & 10 & 4 & 0.64 &  7 & 0.359 & 11 & 0 \\
MOD\_D &  86 & 258 & 10 & 4 & 0.09 & 30 & 0.372 & 18 & 1 \\
\bottomrule
\end{tabular}
\end{table}

\subsection{Working-Day Configuration}
\label{sec:shiftconfig}

For consistency, all instances are evaluated under the same
working-day configuration: 14 working hours per calendar day staffed
by 6 engineers, over a 30-day horizon. The horizon is set generously
so that it does not itself constrain the schedule.

\subsection{Methodology and Metrics}
\label{sec:methodology}

For each instance, fifteen scheduling runs are conducted: one each
of MTS and SPT; ten DEH runs over the $5 \times 2$ grid of weights
$w \in \{0, 0.25, 0.5, 0.75, 1.0\}$ and CST settings
$\gamma \in \{0, 0.5\}$; and three independent runs of CP-SAT under
a 600-second wall-clock limit. Only CP-SAT has solver-internal
variance, so only it is repeated. The $\gamma \in \{0, 0.5\}$ split
lets the effect of CST be recovered as a paired difference per
$(\text{instance}, w)$. The 600-second budget is generous, so any
instance CP-SAT fails to close within it is treated as time-limited.

For each run we record whether a feasible schedule was produced, the
makespan in calendar days, the algorithm-internal runtime in
milliseconds, and, for CP-SAT, its best bound and proved status. The
optimality gap is $(m_{\text{algo}} - m^{*}) / m^{*}$, where $m^{*}$ is
CP-SAT's best feasible makespan on the instance: the proved optimum
where CP-SAT terminates within the budget, the best incumbent at the
cap otherwise. Where CP-SAT does not prove optimality, DEH's gap
against its proved lower bound is also reported as a conservative
cross-check (Section~\ref{sec:per-instance-results}).

A CP-SAT budget sweep characterises the speed--quality trade-off
directly: one additional CP-SAT run per $(\text{instance}, \tau)$ pair
across the grid $\tau \in \{0.1, 1, 10, 60, 600\}$ seconds. The
heuristic is not re-run, since DEH at $(w=0, \gamma=0)$ completes in
under 100 ms on every instance and so already covers every budget at
or above its construction cost.

\subsection{Hypothesis Evaluation}

Throughout, ``best-configured DEH'' denotes the configuration with
the lowest mean gap across the four instances, applied identically to
every instance; here this is $(w, \gamma) = (0, 0)$. This is an
in-sample choice; fixing the configuration before any evaluation runs
is left for the larger instance set noted as future work.

H1 is evaluated by comparing best-configured DEH's mean gap against
those of MTS and SPT; H2 by checking that this mean does not exceed
15\%. The CST ablation ($\gamma = 0.5$ vs $\gamma = 0$ at fixed $w$,
paired per instance) is reported alongside the H1/H2 verdicts.

\subsection{Implementation}

All algorithms are implemented in pure Python and depend only on the
\texttt{ortools} library (CP-SAT, version $\geq 9.8$, $<10$).
Experiments are run on an Apple M1 Pro with 16 GB of RAM under macOS
26.3 with Python 3.13.3; all runtimes are wall-clock medians,
excluding I/O. Source code, the pinned dependency specification, raw
per-instance result CSVs, and aggregation scripts are provided as
supplementary material.

% ============================================================
\section{Results}
\label{sec:results}

\subsection{Per-Instance Results}
\label{sec:per-instance-results}

Table~\ref{tab:per-instance} reports the headline numbers per
instance: MTS, SPT, the lowest-gap DEH configuration on
that instance (with winning $(w, \gamma)$ noted), and the median
CP-SAT result over three independent runs. The full $5 \times 2$
DEH grid is included in the supplementary CSVs.

\begin{table}[ht]
\caption{Per-instance headline results. CP-SAT runtimes are medians
over three runs, which return identical makespans on every instance.
On MOD\_B the CP-SAT entry is the best feasible solution at the
600-second cap, not a proved optimum.}
\label{tab:per-instance}
\centering
\renewcommand{\arraystretch}{1.1}
\setlength{\tabcolsep}{3pt}
\footnotesize
\begin{tabular}{@{}llrrr@{}}
\toprule
Inst.\ & Algorithm                        & $m$ (d) & Gap (\%) & Runtime (ms) \\
\midrule
MOD\_A & MTS                            & 21 & 10.5 & 5    \\
       & SPT                            & 26 & 36.8 & 3    \\
       & DEH $(w{=}0, \gamma{=}0)$      & 20 &  5.3 & 54   \\
       & CP-SAT (med.\ of 3)             & 19 &  0.0 & 4235 \\
\midrule
MOD\_B & MTS                            & 13 & 18.2 & 5      \\
       & SPT                            & 14 & 27.3 & 4      \\
       & DEH $(w{=}0, \gamma{=}0)$      & 12 &  9.1 & 81     \\
       & CP-SAT (med.\ of 3, capped)     & 11 &  0.0 & 600406 \\
\midrule
MOD\_C & MTS                            &  5 &  0.0 & 0    \\
       & SPT                            &  7 & 40.0 & 0    \\
       & DEH $(w{=}0, \gamma{=}0)$      &  5 &  0.0 & 2    \\
       & CP-SAT (med.\ of 3)             &  5 &  0.0 & 76   \\
\midrule
MOD\_D & MTS                            & 11 & 10.0 & 5    \\
       & SPT                            & 12 & 20.0 & 3    \\
       & DEH $(w{=}0, \gamma{=}0)$      & 11 & 10.0 & 22   \\
       & CP-SAT (med.\ of 3)             & 10 &  0.0 & 1059 \\
\bottomrule
\end{tabular}
\end{table}

On every instance the critical-path-only DEH configuration
($w = 0$, $\gamma = 0$) is the lowest-gap heuristic, and SPT gives
the largest gaps. DEH at $w = 0$ strictly beats MTS on two instances
(MOD\_A: 20 vs 21 days; MOD\_B: 12 vs 13 days) and ties on the other
two. CP-SAT proves optimality in seconds on MOD\_A, MOD\_C, and
MOD\_D, but exhausts its 600-second budget on MOD\_B with incumbent
11 and a proved lower bound of 9.

On MOD\_B, where CP-SAT's incumbent (11) is not proved optimal,
DEH's gap against the proved lower bound ($\text{LB} = 9$) is
$(12 - 9)/9 = 33.3\%$. Substituting this lower bound for MOD\_B's
$m^{*}$ raises the cross-instance mean gap from 6.1\% to 12.2\%,
still below the 15\% threshold (H2).

\subsection{Cross-Instance Summary}
\label{sec:cross-instance}

Table~\ref{tab:h1summary} aggregates the results into
hypothesis-evaluation form. The Within-15\% column reports the
fraction of instances whose individual gap is at most 15\%, a
per-instance strengthening of H2's mean constraint.

\begin{table}[ht]
\caption{Cross-instance hypothesis summary. DEH configurations at
intermediate $w$ are omitted for space.}
\label{tab:h1summary}
\centering
\renewcommand{\arraystretch}{1.15}
\setlength{\tabcolsep}{4pt}
\begin{tabular}{@{}lrrr@{}}
\toprule
Configuration                         & Mean gap (\%) & Median gap (\%) & Within 15\% \\
\midrule
MTS                                 &  9.7 & 10.3 & 0.75 \\
SPT                                 & 31.0 & 32.0 & 0.00 \\
DEH $(w{=}0, \gamma{=}0)$           &  6.1 &  7.2 & 1.00 \\
DEH $(w{=}0, \gamma{=}0.5)$         &  7.4 &  9.6 & 1.00 \\
CP-SAT                               &  0.0 &  0.0 & 1.00 \\
\bottomrule
\end{tabular}
\end{table}

\paragraph{H1 (supported)}
The lowest-mean-gap DEH configuration, $w = 0, \gamma = 0$, has a
6.1\% mean gap, below MTS at 9.7\% and SPT at 31.0\%: a 37\% relative
reduction versus the best classical priority rule. DEH at
$(w=0, \gamma=0)$ also never degrades against MTS or SPT on any
individual instance, which H1 does not by itself require.

\paragraph{H2 (supported)}
H2 (mean $\leq 15\%$) is met at 6.1\%. As a stronger sanity check,
every individual instance also lies within the threshold
(5.3\%, 9.1\%, 0\%, 10\%).

\paragraph{Ablation: the constraint-tightness term}
Across the 20 paired comparisons of $\gamma = 0.5$ versus
$\gamma = 0$ at fixed $w$, the split is 3 wins, 11 ties, and 6 losses
for CST: equal-or-better in 14 of 20, but with strict losses
outnumbering strict wins two-to-one. At the recommended $w = 0$,
$\gamma = 0.5$ degrades MOD\_A by one day and ties elsewhere, a
mean-gap delta of $+1.3$ percentage points. Interpretation is
deferred to Section~\ref{sec:discussion}.

\subsection{CP-SAT Time-Budget Sweep}
\label{sec:budget-sweep}

The headline comparison contrasts CP-SAT at its 600-second cap
against the heuristic at $\leq 100$ ms. This is the quality-ceiling
comparison appropriate for the H1 verdicts, not the sub-second regime
motivated in Section~\ref{sec:rqs}. Table~\ref{tab:budget-sweep}
reports the CP-SAT budget sweep.

\begin{table*}[!t]
\caption{CP-SAT makespan (days) across the budget grid
$\tau \in \{0.1, 1, 10, 60, 600\}$ s, versus the heuristic reference
at $\leq 100$ ms. Each cell shows $m$ / status, where status is
\texttt{opt} (CP-SAT proved optimality), \texttt{inc} (feasible
incumbent without proof), or \texttt{none} (no feasible solution
within budget).}
\label{tab:budget-sweep}
\centering
\renewcommand{\arraystretch}{1.15}
\begin{tabular}{@{}lccccccc@{}}
\toprule
       & \multicolumn{5}{c}{CP-SAT at budget $\tau$ (s)}
       & Heuristic ($\leq$0.1 s) & \\
\cmidrule(lr){2-6}
Inst.\ & $\tau{=}0.1$ & $\tau{=}1$ & $\tau{=}10$ & $\tau{=}60$ & $\tau{=}600$
       & DEH $(w{=}0, \gamma{=}0)$ & Best-known $m^{*}$ \\
\midrule
MOD\_A & --/none & 19/inc & 19/opt & 19/opt & 19/opt
       & 20 & 19 (proved) \\
MOD\_B & --/none & 12/inc & 11/inc & 11/inc & 11/inc
       & 12 & 11 (LB 9, unproved) \\
MOD\_C & 5/opt   & 5/opt  & 5/opt  & 5/opt  & 5/opt
       & 5  & 5 (proved) \\
MOD\_D & --/none & 11/inc & 10/opt & 10/opt & 10/opt
       & 11 & 10 (proved) \\
\bottomrule
\end{tabular}
\end{table*}

The sweep separates three regimes. At $\tau = 0.1$ s, CP-SAT returns
no feasible solution on MOD\_A, MOD\_B, and MOD\_D; only MOD\_C is
small enough to be proved optimal. The heuristic returns a feasible
schedule on every instance, so in the sub-second regime it is the
only choice that yields a usable answer on three of the four
instances. At $\tau = 1$ s, CP-SAT matches the heuristic on MOD\_B,
MOD\_C, and MOD\_D and beats it by one day on MOD\_A. At
$\tau \geq 10$ s, CP-SAT reaches its best incumbent everywhere,
beating the heuristic by one day on MOD\_A, MOD\_B, and MOD\_D and
tying on MOD\_C. The heuristic is therefore not uniformly best: it
dominates in the sub-second regime, while CP-SAT weakly dominates
from one second upward when an unproved but high-quality incumbent is
acceptable.

% ============================================================
\section{Discussion}
\label{sec:discussion}

\subsection{Cross-Instance Interpretation}

CP-SAT's behaviour varies sharply with instance structure. On
MOD\_A, MOD\_C, and MOD\_D it proves optimality in under five seconds
despite a roughly five-fold range in task count and DAG depth. On
MOD\_B, which combines a heavy zone footprint (31 zones across 90
tasks) with an exclusive tool and no state activity, it exhausts the
600-second budget; the solver-internal gap is stable across all three
runs rather than reducible by search randomisation. This contrast
motivates the heuristic comparison. Where CP-SAT proves optimality
cheaply, the heuristic gap is the whole question; where it cannot
prove optimality within a feasible budget, an inexpensive heuristic
with a small known gap to the incumbent is operationally useful even
without an exact reference.

The gap shape across $w$ explains why the recommended configuration
sets $w = 0$. At $\gamma = 0$ the per-instance gap is monotone
non-decreasing in $w$ on three of the four instances, so state
continuity alone is a worse ordering criterion than critical-path
priority on this set. CST behaves similarly: it costs one day on
MOD\_A at $w = 0$ and is neutral elsewhere, and its intermediate-$w$
gains (mainly on MOD\_D) do not offset that loss. CST is therefore
retained as a per-instance tunable signal but is neither part of the
recommended configuration nor load-bearing for the headline
comparison.

\subsection{Practical Implications}

The speed gap is decisive for interactive scenario iteration. CP-SAT
saturates the 600-second cap on MOD\_B, so a hundred-scenario sweep
would take well over an hour there, whereas DEH completes each
scenario in under 100\,ms on every instance and finishes the same
sweep in under 10\,s. In this sub-second regime the heuristic is the
only choice that produces a usable schedule on three of the four
instances. Its simplicity brings two further advantages: no external
solver dependency, which suits embedding in a planning tool, and a
small, known gap to CP-SAT that makes it a usable cross-check against
an exact reference.

\subsection{Limitations and Future Work}
\label{sec:limitations}

The strongest limitation is sample size. With $N = 4$, the
cross-instance verdicts (H1, H2) should be read as preliminary
directional evidence rather than statistical generalisation; adequate
coverage of the RCPSP design space requires systematic sampling
across structural parameter axes~\cite{vanhoucke2008measurement},
which four instances do not provide. Each instance is also one module
of one machine, whereas a deployed machine carries several upgradeable
modules; H1, H2, and the CST design ablation can be revisited under
properly powered non-parametric tests, such as a paired Wilcoxon
signed-rank test on the CST deltas, once more instances are available.

MOD\_B is the only instance whose heuristic gap is reported against
an unproved incumbent (cross-checked against the proved lower bound in
Section~\ref{sec:per-instance-results}). Its profile (heavy zone
footprint, an exclusive tool, minimal state activity) motivates
building more instances with similar features to test whether
CP-SAT's difficulty generalises.

The budget sweep uses a coarse five-point grid with a single run per
$(\text{instance}, \tau)$ pair; a finer grid with multi-seed
replication would tighten the curve. The MSA term treats machine
reconfiguration as instantaneous, whereas switching states consumes
time in reality; charging a reconfiguration cost would likely raise
the value of state continuity and shift the optimal $w$ upward, so the
present setup may underestimate MSA. Finally, all instances use one
fixed working-day configuration (14 hours, six staff, 30-day horizon);
a horizon-and-staffing sweep is left for future work.

% ============================================================
\section{Conclusions}
\label{sec:conclusions}

This thesis formalised ASML's upgrade-scheduling problem as a
Resource-Constrained Project Scheduling Problem with five
domain-specific extensions, and evaluated three candidates against
it: a baseline priority-rule greedy (MTS, SPT), an exact CP-SAT
solver, and a domain-enhanced greedy heuristic (DEH).

Both pre-registered hypotheses are supported. The recommended
heuristic configuration, DEH at $w = 0, \gamma = 0$, attains a mean
gap of 6.1\% against the CP-SAT incumbent, strictly below MTS (9.7\%)
and SPT (31.0\%), so H1 holds; this mean is well under the
pre-registered 15\% threshold, so H2 holds; and it is a 37\% relative
reduction versus the best classical priority rule. The CST ablation
is at best conditional: $\gamma = 0.5$ is equal-or-better in 14 of 20
paired comparisons but loses where it matters, degrading the
recommended $w = 0$ mean gap by 1.3 percentage points, so CST is
retained only as an optional per-instance signal.

Returning to the research question, a domain-enhanced greedy
heuristic does close a substantial part of the gap that classical
priority rules leave against CP-SAT, at sub-second cost. The budget
sweep places this in context: below one second CP-SAT returns no
feasible schedule on three of the four instances, so the heuristic is
the only usable option there, whereas from one second upward CP-SAT
matches or beats it on every instance and remains the gold-standard
reference. The heuristic's operational value is therefore to own the
sub-second regime, to remove the external-solver dependency for
embedded use, and to serve as a cheap cross-check against an exact
reference. A broader lesson is that domain-specific ordering signals help only
when the objective rewards them: makespan alone does not, so the
critical path dominates, while an objective that prices machine
reconfiguration would change that calculus.

The verdict is preliminary at $N = 4$, and several directions follow.
The most immediate is scale: a larger, structurally sampled instance
set would move H1 and H2 from directional evidence toward statistical
generalisation, and would support predicting from structural features
alone when the heuristic is safe to trust. The reconfiguration-cost
objective just noted is the natural setting in which to revisit the
machine-state term, which earns nothing under the present
instantaneous-reconfiguration assumption. Finally, deployment spans
several modules per machine and many machines, opening per-machine and
cross-machine scheduling beyond the single-module instances studied
here; finer budget sweeps and multi-seed CP-SAT replication would
tighten the empirical picture throughout.

% ============================================================
\begin{thebibliography}{00}

\bibitem{blazewicz1983scheduling}
J.~Blazewicz, J.~K.~Lenstra, and A.~H.~G.~Rinnooy~Kan,
``Scheduling subject to resource constraints: classification and
complexity,''
\textit{Discrete Applied Mathematics}, vol.~5, no.~1, pp.~11--24, 1983.

\bibitem{kolisch1997psplib}
R.~Kolisch and A.~Sprecher,
``PSPLIB: a project scheduling problem library,''
\textit{European Journal of Operational Research}, vol.~96, no.~1,
pp.~205--216, 1997.

\bibitem{laborie2018ibm}
P.~Laborie, J.~Rogerie, P.~Shaw, and P.~Vil\'im,
``IBM ILOG CP optimizer for scheduling,''
\textit{Constraints}, vol.~23, no.~2, pp.~210--250, 2018.

\bibitem{Baptiste2001}
P.~Baptiste, C.~Le~Pape, and W.~Nuijten,
\textit{Constraint-Based Scheduling: Applying Constraint Programming
to Scheduling Problems}.
Springer, 2001.

\bibitem{hartmann2010survey}
S.~Hartmann and D.~Briskorn,
``A survey of variants and extensions of the resource-constrained
project scheduling problem,''
\textit{European Journal of Operational Research}, vol.~207, no.~1,
pp.~1--14, 2010.

\bibitem{kolisch1996efficient}
R.~Kolisch,
``Efficient priority rules for the resource-constrained project
scheduling problem,''
\textit{Journal of Operations Management}, vol.~14, no.~3,
pp.~179--192, 1996.

\bibitem{ortools}
Google OR-Tools,
``OR-Tools: open source software for combinatorial optimization,''
Google Developers, 2023. [Online].
Available: https://developers.google.com/optimization

\bibitem{brucker1999rcpsp}
P.~Brucker, A.~Drexl, R.~M\"ohring, K.~Neumann, and E.~Pesch,
``Resource-constrained project scheduling: notation, classification,
models, and methods,''
\textit{European Journal of Operational Research}, vol.~112, no.~1,
pp.~3--41, 1999.

\bibitem{kolisch2006experimental}
R.~Kolisch and S.~Hartmann,
``Experimental investigation of heuristics for resource-constrained
project scheduling: an update,''
\textit{European Journal of Operational Research}, vol.~174, no.~1,
pp.~23--37, 2006.

\bibitem{stinson1978multiple}
J.~P.~Stinson, E.~W.~Davis, and B.~M.~Khumawala,
``Multiple resource-constrained scheduling using branch and bound,''
\textit{AIIE Transactions}, vol.~10, no.~3, pp.~252--259, 1978.

\bibitem{vanhoucke2008measurement}
M.~Vanhoucke, J.~Coelho, D.~Debels, B.~Maenhout, and L.~V.~Tavares,
``An evaluation of the adequacy of project network generators with
systematically sampled networks,''
\textit{European Journal of Operational Research}, vol.~187, no.~2,
pp.~511--524, 2008.

\bibitem{laborie2008complete}
P.~Laborie and J.~Rogerie,
``Reasoning with conditional time-intervals,''
in \textit{Proc.\ 21st Int.\ FLAIRS Conference}, 2008, pp.~555--560.

\bibitem{vilim2015fastcumulative}
P.~Vil\'im,
``Timetable edge finding filtering algorithm for discrete cumulative
resources,''
in \textit{Integration of AI and OR Techniques in Constraint
Programming (CPAIOR)}, 2015, pp.~230--245.

\bibitem{demeulemeester2002project}
E.~Demeulemeester and W.~Herroelen,
\textit{Project Scheduling: A Research Handbook}.
Kluwer Academic, 2002.

\bibitem{megow2011turnaround}
N.~Megow, R.~H.~M\"ohring, and J.~Schulz,
``Decision support and optimization in shutdown and turnaround
scheduling,''
\textit{INFORMS Journal on Computing}, vol.~23, no.~2, pp.~189--204,
2011.

\bibitem{monch2011semiconductor}
L.~M\"onch, J.~W.~Fowler, S.~Dauz\`ere-P\'er\`es, S.~J.~Mason, and
O.~Rose,
``A survey of problems, solution techniques, and future challenges in
scheduling semiconductor manufacturing operations,''
\textit{Journal of Scheduling}, vol.~14, no.~6, pp.~583--599, 2011.

\end{thebibliography}

\end{document}
